\section{Conclusion}

This paper presented Marine Race Arena, a configurable simulation-oriented benchmark layer for underwater multi-robot racing. The proposed system adapts the repeatable, gate-based evaluation logic of autonomous racing to underwater vehicles while adding marine-specific elements such as acoustic guidance, depth constraints, currents, obstacles, and referee-controlled scoring. The central contribution is not a new simulator or a new controller, but a modular evaluation architecture in which controllers can be replaced while arena rules, task definitions, logs, and scoring remain fixed.

The current draft establishes the article structure, related work positioning, system architecture, environment design, evaluation protocol, and baseline status. The next required step is quantitative validation: repeated benchmark runs for acoustic-only and rule-based vision-assisted baselines on clean-gate, obstacle-gate, and current-gate tasks. Once those results are available, Marine Race Arena can be positioned as a credible first step toward reproducible underwater racing benchmarks.
